% TOP_PROBLEM: {"id": 2302004, "title": "Research Problems in Function Theory — Problem 2.4", "category_id": 8, "category": {"id": 8, "name": "analysis", "display_name": "Analysis", "description": "Limits, continuity, calculus, and function theory.", "slug": "analysis", "order_index": 8, "created_at": "2026-07-31T15:26:25.671Z"}, "status": "open"}
% FIRST_POSTED: null
% ENTRY_KIND: statement_only
% Imported from: batch10.tex; UnsolvedMath ID 2302004
\documentclass[11pt]{article}
\usepackage[margin=25mm]{geometry}
\usepackage{fontspec}
\setmainfont{Latin Modern Roman}
\usepackage{amsmath,amssymb,mathtools}
\usepackage{unicode-math}
\setmathfont{Latin Modern Math}
\usepackage{xurl,hyperref}
\hypersetup{colorlinks=true,urlcolor=blue}
\setcounter{secnumdepth}{0}
\setlength{\parindent}{0pt}
\setlength{\parskip}{5pt}
\setlength{\emergencystretch}{3em}
\newcommand{\sref}[2]{\hyperlink{source:#1:#2}{[S#2]}}
\newcommand{\cataloguescope}[1]{\paragraph{Recorded target.}#1}
\begin{document}
% UnsolvedMath ID 2302004; source catalogue code AMR-022-2004
\setcounter{section}{2302003}
\section{Research Problems in Function Theory — Problem 2.4}\label{2302004}
{\small\sffamily UnsolvedMath ID 2302004 · Analysis}
\subsection{Definitions and mathematical statement}

\paragraph{Shared notation.}$\mathbb N=\{1,2,\ldots\}$, and $\mathbb Z,\mathbb Q,\mathbb R,\mathbb C$ denote the integers, rationals, reals and complex numbers. Any other conventions are specified locally.


For a nonconstant meromorphic function $f$ on $\mathbb C$, let $n(r,a,f)$ count the solutions of $f(z)=a$ in $|z|\le r$, with multiplicity; $a=\infty$ counts poles. Let $\bar n$ count each point once. The integrated count is
\[
 N(r,a,f)=n(0,a,f)\log r+\sum_{0<|z|\le r,\ f(z)=a}\log(r/|z|),
\]
where the sum repeats each point according to multiplicity; $\bar N$ has the same formula without repetitions. Write $\log^+x=\max(\log x,0)$ and
\[
 m(r,\infty,f)=\frac1{2\pi}\int_0^{2\pi}\log^+|f(re^{it})|\,dt,\quad
 m(r,a,f)=\frac1{2\pi}\int_0^{2\pi}\log^+\frac1{|f(re^{it})-a|}\,dt\quad(a\in\mathbb C).
\]
The Nevanlinna characteristic is $T(r,f)=m(r,\infty,f)+N(r,\infty,f)$. Its order and lower order are
\[
 \rho(f)=\limsup_{r\to\infty}\frac{\log T(r,f)}{\log r},\qquad
 \lambda(f)=\liminf_{r\to\infty}\frac{\log T(r,f)}{\log r}.
\]
The deficiency is $\delta(a,f)=1-\limsup_{r\to\infty}N(r,a,f)/T(r,f)$, and the reduced-count index is $\Theta(a,f)=1-\limsup_{r\to\infty}\bar N(r,a,f)/T(r,f)$. A value is deficient if $\delta(a,f)>0$. An entire function is holomorphic on all of $\mathbb C$; for an entire function, $M(r,f)=\max_{|z|=r}|f(z)|$.

For angles $\theta$ modulo $2\pi$ and $0<\varepsilon<\pi$, put
\[
 D(\theta,\varepsilon)=\{z\ne0:\min_{j\in\mathbb Z}|\arg z-\theta+2\pi j|<\varepsilon\}.
\]
Saying that $f$ assumes $a$ infinitely often in this sector means that the equation $f(z)=a$ has infinitely many distinct solutions there, with poles interpreted as $\infty$-points.

For a nonconstant meromorphic $f$, define
\[
 E_f(\theta)=\{a\in\widehat{\mathbb C}:\text{for some }\varepsilon>0,
                  f(z)=a\text{ has only finitely many solutions in }D(\theta,\varepsilon)\}.
\]
The ray of angle $\theta$ is a Julia line when $|E_f(\theta)|\le2$. For an entire function $\infty\in E_f(\theta)$, so a Julia line has at most one finite exceptional value.
\paragraph{Statement recorded in the supplied source.}
Determine the restrictions on the exceptional-value sets at different Julia lines of a meromorphic function. In particular, can an entire function have Julia lines of distinct angles $\theta_a,\theta_b$ with distinct finite exceptional values $a\in E_f(\theta_a)$ and $b\in E_f(\theta_b)$? \sref{2302004}{2}
\subsection{Short English statement}
Compare exceptional values at different Julia directions and ask whether two different finite exceptions can occur for one entire function. The source reports examples with different exceptions at each of finitely many Julia lines. \sref{2302004}{2}
\subsection{Sources}
\begin{description}
\item[\hypertarget{source:2302004:1}{S1}] ULAM.AI and the UnsolvedMath Contributors, UnsolvedMath: A Curated Collection of Open Mathematics Problems, record 2302004 (AMR-022-2004). CC BY 4.0. \url{https://huggingface.co/datasets/ulamai/UnsolvedMath}.
\item[\hypertarget{source:2302004:2}{S2}] Walter K. Hayman and Eleanor F. Lingham, Research Problems in Function Theory, Fiftieth Anniversary Edition (2018), arXiv:1809.07200. Problem 2.4, its full Update 2.4, and chapter notation; original author TeX. \url{https://arxiv.org/abs/1809.07200}.
\end{description}
\label{end:2302004}
\end{document}
